\section[Demonstrations]{\secstyle{Demonstrations:}}

\section[Exercices]{\secstyle{Exercices:}}

\subsection[Exercice 25]{\subsecstyle{Exercice 25:}}
\subsubsection*{\subsubsecstyle{(3a) Montrer que \(f \circ g^-1 = g^-1 \circ f \):}}

\begin{proof}[\unskip\nopunct]
   Partant de l'hypothèse \(g = f \circ f \), on a:
   \begin{flalign*}
      & & f \circ f \circ f &= f \circ g  \shorteqnote{(Composition par f par la gauche)}\\
      & & g \circ f & = f \circ g \shorteqnote{(Par définition de g)}\\
      & & g \circ f \circ g^{-1} & = f \circ g \circ g^{-1} \shorteqnote{(Composition par \(g^{-1}\) par la droite)}\\
      & & g \circ f \circ g^{-1} & = f \circ Id_E \\
      & & g^{-1} \circ g \circ f \circ g^{-1} & = g^{-1} \circ f \shorteqnote{(Composition par \(g^{-1}\) par la gauche)}\\
      & & Id_E \circ f \circ g^{-1} & = g^{-1} \circ f \\
      & & f \circ g^{-1} & = g^{-1} \circ f \\
   \end{flalign*}
\end{proof}

\subsubsection*{\subsubsecstyle{(3b) Une expression de \(f \circ g^{-1} \circ f\):}}
On a donc \(f \circ g^{-1} = g^{-1} \circ f\), puis:

\begin{flalign*}
   & &f \circ g^{-1} \circ f &= g^{-1} \circ f \circ f \shorteqnote{(Composition par f par la droite)}\\
   & &f \circ g^{-1} \circ f &= g^{-1} \circ g\\
   & &f \circ g^{-1} \circ f &= Id_E
\end{flalign*}

\subsubsection*{\subsubsecstyle{(3c) En déduire que \(f\) est bijective et expliciter \(f^{-1}\) en fonction de \(f\) et de \(g^{-1}\):}}
\begin{proof}[\unskip\nopunct]
   Soit \(h_1 := g^{-1} \circ f\) et \(h_2 := f \circ g^{-1}\), on a montré dans la première question que \(h_1 = h_2\), et dans la seconde question que \(f \circ h_1 = Id_E\), mais aussi que \(h_2 \circ f = Id_E\).\<
   
   Donc on a bien trouvé une application \(h\) telle que:
   \[
      f \circ h = h \circ f = Id_E
   \]
   On en conclut donc que f est bijective.
\end{proof}
